\chapter{Tweede Deel: Dialectiek van het esthetisch oordeelsvermogen}

\section{.}

Een oordeelskracht, die dialectisch moet zijn, moet allereerst rationalistisch zijn, d.w.z. haar oordelen moeten aanspraak maken op algemene geldigheid, en wel \textit{a priori},\footnote{
Een rationalistisch oordeel (\textit{iudicium ratiocinans}) is elk oordeel dat zichzelf als universeel verklaart, want in die mate kan het dienen als grote premisse in een redenering van de rede. Daartegenover kan alleen een oordeel dat wordt gedacht als de conclusie van een redenering van de rede, en bijgevolg \textit{a priori} gefundeerd is, een oordeel van de rede (\textit{iudicium ratiocinatum}) worden genoemd.} want de dialectiek bestaat in de tegenstrijdigheid van dergelijke oordelen. Derhalve is de onverenigbaarheid van esthetische oordelen van zintuig (over het aangename en het onaangename) niet dialectisch. Zelfs de strijd tussen smaakoordelen, voor zover ieder zich slechts op zijn eigen smaak beroept, vormt geen dialectiek van de smaak, omdat niemand er ook maar aan denkt, zijn eigen oordeel tot een universele regel te maken. Er blijft dus geen ander begrip van een dialectiek over, dat op de smaak van toepassing zou kunnen zijn, dan dat van een dialectiek van de \textbf{kritiek} van de smaak (niet van de smaak zelf) met betrekking tot haar \textbf{principes}: namelijk, waar wederzijds strijdige begrippen van de grond van de mogelijkheid van smaakoordelen natuurlijk en onvermijdelijk naar voren komen. Een transcendentale kritiek van de smaak zal dus een deel bevatten, dat de naam van een dialectiek van de esthetische oordeelskracht kan dragen, alleen als er een antinomie van de principes van dit vermogen is, die haar wetmatigheid en dus ook haar innerlijke mogelijkheid twijfelachtig maakt.

\begin{mdframed}[leftmargin=0pt,rightmargin=0pt,backgroundcolor=blue!6,linecolor=white,innertopmargin=6pt,innerbottommargin=6pt]

Een oordeel dat dialectisch wil zijn – dat wil zeggen, een oordeel dat tot tegenstrijdigheden kan leiden – moet allereerst redelijk zijn. Dat betekent dat het aanspraak maakt op universele geldigheid, en wel zonder afhankelijk te zijn van ervaring (a priori). Dialectiek ontstaat immers wanneer dergelijke oordelen met elkaar in strijd zijn.

Daarom kan de onverenigbaarheid van zintuiglijke oordelen (zoals ``dit is aangenaam'' of ``dit is onaangenaam'') nooit echt dialectisch zijn. Die oordelen zijn immers subjectief en maken geen aanspraak op algemene geldigheid. Ook de strijd tussen smaakoordelen, waarbij ieder zich beroept op zijn eigen smaak, vormt geen echte dialectiek. Niemand probeert immers zijn persoonlijke smaak op te leggen als een universele regel.

Er blijft dus maar één mogelijkheid over voor een dialectiek die op smaak betrekking heeft: niet een dialectiek van de smaak zelf, maar een dialectiek van de kritiek op de smaak. Dit doet zich voor wanneer er tegenstrijdige principes naar voren komen over de grond waarop smaakoordelen eigenlijk mogelijk zijn. Een transcendentale kritiek van de smaak zal dus alleen een dialectisch deel bevatten als er sprake is van een antinomie – een schijnbare tegenstrijdigheid tussen principes – die de wetmatigheid en dus de innerlijke mogelijkheid van smaakoordelen in twijfel trekt.

\bigskip
In §55 komt Kants diepste zorg over de esthetische oordeelskracht naar voren: de vraag of schoonheid en smaak wel objectief gefundeerd kunnen zijn, of dat ze uiteindelijk slechts subjectieve voorkeuren blijven. Deze paragraaf is geen losstaand stuk, maar een cruciaal moment in zijn hele \textit{Kritik der Urteilskraft}. Om dat te begrijpen, moeten we kijken naar de structuur en functie van het gehele werk, en naar de plaats die deze paragraaf daarin inneemt.

Kant deelt zijn werk op in twee grote delen, die elk een fundamenteel vermogen van de menselijke geest onderzoeken: de esthetische oordeelskracht (schoonheid en het verhevene) en de teleologische oordeelskracht (doelmatigheid in de natuur). Beide delen volgen eenzelfde logische structuur, die Kant ook in zijn eerdere Kritieken hanteert: analytiek, dialectiek en methodologie. Deze indeling is geen willekeurige keuze, maar een spiegel van hoe de menselijke geest werkt, hoe wij betekenis geven aan de wereld om ons heen.
Het eerste deel, \textit{Kritiek van de esthetische oordeelskracht}, onderzoekt hoe wij schoonheid en het verhevene ervaren; niet als louter subjectieve gevoelens, maar als universele ervaringen die ons verbinden met anderen en met de wereld.

In het eerste boek daarbinnen, \textit{Analytiek van het schone}, ontleedt Kant hoe wij schoonheid ervaren. Hij laat zien dat een esthetisch oordeel (zoals ``dit landschap is mooi'') geen louter persoonlijke mening is, maar een oordeel dat aanspraak maakt op universele instemming. Schoonheid is geen kwestie van smaak, maar van een vrije harmonie tussen verbeelding en verstand. Dit boek is de kern van Kants esthetica: het toont aan dat schoonheid niet willekeurig is, maar berust op structuren in ons denkvermogen.

In het tweede boek, \textit{Analytiek van het verhevene}, onderzoekt Kant het verhevene: de ervaring van grootsheid of overweldigende kracht (zoals een storm op zee of een reusachtig gebergte), die ons niet alleen raakt, maar ook uitdaagt om ons eigen vermogen te overstijgen. Het verhevene is geen schoonheid, maar een ervaring van onze eigen vrijheid in het aangezicht van het oneindige. Dit boek laat zien dat esthetica niet alleen over harmonie gaat, maar ook over conflict en overwinning – over hoe wij ons verhouden tot datgene wat ons overstijgt.

De \textit{Dialectiek van de esthetische oordeelskracht}, (waaronder §55 valt), is het kritische deel van Kants esthetica. Hier onderzoekt hij of er tegenstrijdigheden (antinomieën) zijn in de principes waarop onze esthetische oordelen berusten. §55 is hierbij cruciaal: Kant vraagt zich af of er wel een echte dialectiek mogelijk is in de smaak. Zijn conclusie is dat er geen echte tegenstrijdigheid is in smaakoordelen zelf (want niemand dwingt zijn smaak op als universele regel), maar wel in de kritiek op de smaak – dat wil zeggen, in de vraag of smaak wel objectief gefundeerd kan zijn. Dit is een diepgaande reflectie op de grenzen van esthetische oordelen: kunnen wij wel echt zeggen dat schoonheid universeel is, of is het uiteindelijk toch subjectief? Deze dialectiek is geen afwijzing van schoonheid, maar een onderzoek naar de voorwaarden waarop esthetische oordelen mogelijk zijn.
De functie van dit deel is dus kritisch: het toont de grenzen van esthetische oordelen en bereidt de weg voor naar de teleologie, de vraag of er een doelmatige orde in de natuur is, die onze esthetische ervaringen kan rechtvaardigen.

\end{mdframed}

\section{Voorstelling van de antinomie van de smaak.}

Het eerste gemeenplaats van de smaak luidt: \textbf{Ieder heeft zijn eigen smaak}. Dit is het uitvluchtmiddel van wie geen smaak heeft en zich tegen kritiek wil verdedigen. Het betekent niets anders dan dat de bepalende grond van dit oordeel louter subjectief is (genot of pijn), en dat het oordeel geen recht heeft op de noodzakelijke instemming van anderen.

Het tweede gemeenplaats, dat zelfs door hen wordt gebruikt die aan smaakoordelen het recht toekennen om geldig te zijn voor iedereen, is: \textbf{Over smaak valt niet te twisten}. Dit betekent dat de bepalende grond van een smaakoordeel weliswaar objectief *kan* zijn, maar niet tot bepaalde begrippen kan worden teruggebracht; bijgevolg kan er niets over het oordeel zelf door middel van bewijzen worden \textbf{beslist}, hoewel het zeker mogelijk en rechtmatig is erover te \textbf{redeneren}. Want \textbf{redeneren} en \textbf{twisten} lijken weliswaar op elkaar, doordat ze beide proberen om door wederzijdse tegenstelling overeenstemming in oordelen teweeg te brengen, maar ze verschillen doordat het laatste dit hoopt te bereiken volgens bepaalde begrippen als gronden van bewijzen, en dus \textbf{objectieve begrippen} als gronden van het oordeel veronderstelt. Waar dit onmogelijk wordt geacht, wordt ook twisten onmogelijk geacht.

Het is gemakkelijk in te zien dat tussen deze twee gemeenplaatsen een stelling ontbreekt, die weliswaar geen spreekwoord is dat algemeen in omloop is, maar waar iedereen toch enigszins een gevoel voor heeft: \textbf{Over smaak valt wel te redeneren} (maar niet te twisten). Maar deze stelling sluit het tegendeel in van de eerste stelling hierboven. Want waar het mogelijk wordt geacht om te redeneren, moet er hoop zijn op wederzijdse overeenstemming; men moet dus kunnen rekenen op gronden voor het oordeel die niet louter private geldigheid hebben en dus niet louter subjectief zijn, wat toch geheel in strijd is met het fundamentele principe: \textbf{Ieder heeft zijn eigen smaak}.

Er bestaat dus de volgende antinomie met betrekking tot het principe van de smaak:

\begin{enumerate}
\item \textbf{These}. Het smaakoordeel berust niet op begrippen, want anders zou erover getwist kunnen worden (door middel van bewijzen beslist).
\item \textbf{Antithese}. Het smaakoordeel berust wel op begrippen, want anders zou er, ondanks de verscheidenheid ervan, niet eens over kunnen worden geredeneerd (aanspraak kunnen worden gemaakt op de noodzakelijke instemming van anderen bij dit oordeel).
\end{enumerate}

\begin{mdframed}[leftmargin=0pt,rightmargin=0pt,backgroundcolor=blue!6,linecolor=white,innertopmargin=6pt,innerbottommargin=6pt]

Wanneer het over smaak gaat, hoor je vaak twee uitspraken die als vanzelfsprekend worden aangenomen, alsof ze onbetwistbare waarheden zijn. De eerste is het bekende ``\textit{Ieder heeft zijn eigen smaak}'', een uitvlucht die vooral wordt gebruikt door wie geen verweer heeft tegen kritiek op hun oordeel. Wat hiermee eigenlijk wordt gezegd, is dat smaak louter een kwestie is van persoonlijke voorkeur, van wat iemand aangenaam of onaangenaam vindt. Het is een bewering die elk smaakoordeel reduceert tot een subjectief gevoel, zonder enige aanspraak op universele geldigheid. Als smaak echt niets meer is dan individuele gril, dan heeft het geen zin om erover in discussie te gaan, want wie kan nu bepalen wat een ander mooi vindt? Maar deze opvatting ontkent precies wat wij allemaal diep vanbinnen voelen: dat schoonheid niet alleen een persoonlijke mening is, maar iets dat ook anderen zou moeten aanspreken. Wanneer wij zeggen ``\textit{dit schilderij is mooi}'', dan bedoelen wij niet alleen ``\textit{ik vind het mooi}'', maar ook ``\textit{jij zou het eigenlijk ook mooi moeten vinden}''. Dat is geen arrogantie, maar een wezenlijk kenmerk van elk esthetisch oordeel: het doet een beroep op instemming, alsof schoonheid niet alleen *mijn* ervaring is, maar iets dat voor iedereen geldt.

De tweede gemeenplaats is niet minder problematisch: ``\textit{Over smaak valt niet te twisten.}'' Ook wie erkent dat smaakoordelen een zekere objectiviteit kunnen hebben, grijpt naar deze uitdrukking om aan te geven dat je niet met harde bewijzen kunt aantonen wat mooi is. En dat klopt: je kunt niet, zoals in de wiskunde, met logische redeneringen bewijzen dat een landschap of een symfonie schoonheid bezit. Toch is dit geen reden om te zeggen dat er helemaal niet over smaak gesproken kan worden. Er is een wezenlijk verschil tussen twisten en redeneren. Twisten doe je wanneer je met vaste begrippen en onweerlegbare bewijzen een punt probeert te maken, zoals in een wetenschappelijk debat. Maar redeneren, een gesprek voeren waarin je probeert om elkaars standpunt te begrijpen, om tot overeenstemming te komen, dat kan wel over smaak. Sterker nog: het feit dat wij wel over smaak kunnen praten, dat wij elkaar kunnen overtuigen, dat wij soms zelfs tot een gedeelde mening komen, bewijst dat smaak niet louter subjectief is. Als smaak echt alleen maar een kwestie was van persoonlijke voorkeur, dan zou elk gesprek erover zinloos zijn. Toch voelen wij allemaal aan dat dat niet zo is. Wanneer wij met elkaar discussiëren over een gedicht, een gebouw of een muziekstuk, dan gaan wij ervan uit dat er redenen zijn voor ons oordeel, redenen die niet alleen voor onszelf gelden, maar die ook anderen kunnen overtuigen.

Hier ontstaat nu een schijnbare tegenstrijdigheid, een antinomie, die de kern raakt van hoe wij over smaak denken. Aan de ene kant lijkt het alsof smaakoordelen niet berusten op duidelijke begrippen. Als dat wel zo was, dan zouden wij erover kunnen twisten, net als over een wiskundige stelling of een wetenschappelijke hypothese. We kunnen immers niet met logische argumenten aantonen dat een zonondergang mooier is dan een regenboog. Maar aan de andere kant *moeten* smaakoordelen wel op iets berusten dat meer is dan louter persoonlijke smaak, want anders zou het onmogelijk zijn om erover te redeneren, om elkaars mening te beïnvloeden, om tot een gedeelde ervaring van schoonheid te komen. Als smaak echt alleen maar subjectief was, dan zou elk gesprek erover niets meer zijn dan een uitwisseling van onverenigbare meningen, zonder enige hoop op overeenstemming. Maar dat is niet hoe wij smaak ervaren. Wij verwachten dat anderen onze oordelen kunnen delen, wij zoeken naar woorden om onze ervaring uit te drukken, wij hopen dat anderen hetzelfde zullen voelen. Dat betekent dat er wel degelijk een soort objectiviteit in het spel moet zijn; niet de objectiviteit van een wetenschappelijk feit, maar die van een ervaring die universeel mededeelbaar is.

Dus staan wij voor een raadsel: smaakoordelen berusten niet op vaste begrippen die wij kunnen bewijzen, en toch maken zij aanspraak op een instemming die verder gaat dan onze persoonlijke voorkeur. Ze zijn subjectief, maar niet louter subjectief; ze zijn objectief, maar niet op de manier van een wetmatigheid. Dit is de spanning waarin elk esthetisch oordeel zich bevindt, en die spanning is geen zwakte, maar juist de kracht van smaak. Want het is precies deze dubbelheid die ons dwingt om niet alleen te voelen, maar ook te denken over schoonheid, om niet alleen onze eigen reactie te volgen, maar ook te zoeken naar wat anderen ervaren. Schoonheid is geen kwestie van louter gevoel, maar ook niet van louter verstand. Het is een ervaring die ons zowel als individuen als als gemeenschap raakt, die ons uitnodigt om te delen wat wij zien en te zoeken naar wat ons verbindt.

\bigskip


\epigraph{\textit{``Schoonheid is geen kwestie van \textquoteleft dit is mooi\textquoteright\ of \textquoteleft dit is niet mooi\textquoteright, maar van \textquoteleft Laten wij samen kijken, voelen en denken.\textquoteright''}}{--- Impliciete les uit Kants esthetica}

§56 van Kants \textit{Kritik der Urteilskracht} onthult de \textbf{diepste spanning} in ons denken over schoonheid --- een spanning die niet alleen filosofisch, maar ook \textbf{menselijk en spiritueel} is. Deze paragraaf toont aan dat schoonheid voor Kant geen oppervlakkige kwestie is, maar een \textbf{fundamentele manier waarop wij betekenis geven aan de wereld}. Voor vrijmetselaren resoneren deze inzichten bijzonder sterk, omdat zij de essentie raken van wat het betekent om \textbf{te zoeken, te delen en te verbinden}.

\subsection*{De antinomie als spiegel van de menselijke geest}

Kant stelt ons voor een schijnbaar onoplosbare, maar vruchtbare tegenstrijdigheid:
\begin{quote}
\emph{Ieder heeft zijn eigen smaak}, maar tegelijkertijd \emph{verwachten wij instemming} wanneer wij iets mooi noemen. Schoonheid is subjectief, maar maakt aanspraak op universaliteit; zij is niet te bewijzen, maar wel te bespreken.
\end{quote}
Deze spanning is geen fout in ons denken, maar juist het teken van iets diepers: schoonheid is geen zaak van \textbf{weten}, maar van \textbf{ervaren en delen}. Voor vrijmetselaren is dit herkenbaar in het ritueel, dat zowel \textbf{persoonlijk} als \textbf{gemeenschappelijk} is. Net zoals een symbool voor ieder individu een unieke betekenis heeft, en toch allen verbindt, zo is schoonheid een ervaring die \textbf{tussen ons en de wereld} bestaat --- iets dat wij \emph{samen} moeten ontdekken.

\subsection*{Schoonheid als oefening in vrijheid en verbondenheid}

De antinomie leert ons dat:
\begin{itemize}
    \item \textbf{Vrijheid}: Schoonheid dwingt niet. Niemand \emph{moet} iets mooi vinden, maar de ervaring nodigt uit tot \textbf{gedeelde reflectie}.
    \item \textbf{Verbondenheid}: Wanneer wij over smaak redeneren, zoeken wij niet naar bewijzen, maar naar \textbf{gemeenschappelijke grond} --- precies zoals het maçonnieke ideaal van broederschap niet eist dat wij hetzelfde denken, maar dat wij \textbf{samen zoeken}.
\end{itemize}
Kant toont aan dat ware schoonheid niet bestaat in eenduidigheid, maar in dialoog. Dit is de kern van het logeleven: \textbf{geen dogma's, maar gedeelde zoektocht}.

\subsection*{De maçonnieke dimensie: Symbolen als esthetische oordelen}

Maçonnieke symbolen (zoals het vierkant, de passer, of het licht) functioneren net als smaakoordelen:
\begin{itemize}
    \item Zij hebben geen vaste betekenis, maar nodigen uit tot \textbf{persoonlijke interpretatie}.
    \item Toch verwachten wij --- niet als eis, maar als hoop --- dat anderen er ook iets in herkennen.
\end{itemize}
Symbolen werken omdat zij \textbf{open} zijn: zij verbinden zonder voor te schrijven. Dit is de kracht van het maçonnieke pad: \textbf{zoeken zonder te eisen, delen zonder op te leggen}.

\subsection*{Schoonheid als brug tussen individu en gemeenschap}

Kants antinomie is de kern van zijn betoog, omdat zij schoonheid presenteert als \textbf{de brug tussen individu en gemeenschap}. Schoonheid is geen bezit, maar een \textbf{activiteit} --- een manier van \textbf{onderweg zijn}, vrij en verbonden. Dit is de essentie van de vrijmetselarij: ware wijsheid begint niet met vaste antwoorden, maar met \textbf{open vragen en gedeelde ervaring}.

\subsection*{Conclusie: Schoonheid als weg naar wijsheid}

§56 is geen abstracte filosofie, maar een \textbf{gids voor het leven}:
\begin{quote}
Ware kennis komt niet van dogma's, maar van \textbf{samen zoeken, luisteren en groeien}.
\end{quote}
Voor vrijmetselaren is dit de les die in elk ritueel doorklinkt: schoonheid, net als wijsheid, is geen bestemming, maar een \textbf{manier van onderweg zijn} --- vrij en verbonden, individueel en gemeenschappelijk.

\end{mdframed}

\section{Oplossing van de antinomie van de smaak.}

Er is geen andere mogelijkheid om de tegenstrijdigheid tussen deze twee principes, die ten grondslag liggen aan elk smaakoordeel (die niets anders zijn dan de twee bijzonderheden van het smaakoordeel, zoals uiteengezet in de Analytiek), op te heffen, dan door aan te tonen dat het begrip waaraan het voorwerp in dit soort oordeel wordt gerelateerd, niet in dezelfde zin wordt genomen in de twee uitspraken van de esthetische oordeelskracht; dat deze dubbele betekenis of dit gezichtspunt in het oordelen noodzakelijk is in ons transcendentale oordeelsvermogen; en dat de schijn die ontstaat door de verwarring van het ene met het andere een natuurlijke illusie is, die onvermijdelijk is.

Het smaakoordeel moet gerelateerd zijn aan enig concept, want anders zou het geen aanspraak kunnen maken op noodzakelijke geldigheid voor iedereen. Maar het hoeft daarom niet \textbf{uit} een concept afleidbaar te zijn, omdat een concept óf determineerbaar kan zijn, óf in zichzelf onbepaald en ook onbepaalbaar. Het concept van de rede, dat bepaalbaar is door middel van predicaten van de zintuiglijke intuïtie die ermee kunnen overeenstemmen, behoort tot het eerste soort; het transcendentaal concept van de rede van het bovenzintuiglijke, dat de grond vormt van alle intuïtie en dus theoretisch niet verder kan worden bepaald, behoort tot het tweede soort.

Het smaakoordeel heeft weliswaar betrekking op voorwerpen der zintuigen, maar niet om een \textbf{begrip} ervan voor het verstand te bepalen, want het is geen kennisoordeel. Het is dus, als een intuïtieve, enkelvoudige voorstelling, gerelateerd aan het gevoel van genot, slechts een particulier oordeel, en in die mate is zijn geldigheid beperkt tot het oordelende individu alleen: het voorwerp is \textbf{voor mij} een voorwerp van bevrediging, voor anderen kan het anders zijn, - ieder heeft zijn eigen smaak.

Toch bevat het smaakoordeel ongetwijfeld een uitgebreide relatie van de voorstelling van het voorwerp (en tegelijkertijd van het subject), waarop wij een uitbreiding van dit soort oordeel als noodzakelijk voor iedereen baseren, die dus op enig concept moet berusten, maar op een concept dat \textbf{niet} door intuïtie kan worden bepaald, waardoor niets kan worden gekend en dat dus ook tot \textbf{geen bewijs} voor het smaakoordeel leidt. Een dergelijk concept is echter het louter zuivere rationele concept van het bovenzintuiglijke, dat het voorwerp (en ook het oordelende subject) als een voorwerp der zintuigen, dus als verschijning, grondvest. Want als men zo'n gezichtspunt niet zou aannemen, dan zou de aanspraak van het smaakoordeel op universele geldigheid niet te redden zijn; als het concept waarop het berust slechts een verward begrip van het verstand was, bijvoorbeeld van volmaaktheid, waaraan men overeenkomstig de zintuiglijke intuïtie van het schone zou kunnen toekennen, dan zou het mogelijk zijn, althans in zichzelf, het smaakoordeel op bewijzen te gronden, wat in strijd is met de these.

Maar nu verdwijnt alle tegenstrijdigheid als ik zeg dat het smaakoordeel berust op een concept (van een algemene grond voor de subjectieve doelmatigheid van de natuur voor de oordeelskracht), waaruit echter niets over het voorwerp kan worden gekend en bewezen, omdat het in zichzelf onbepaalbaar en ongeschikt is voor kennis; maar dat het tegelijkertijd door middel van dit concept geldigheid voor iedereen verkrijgt (in elk geval als een enkelvoudig oordeel dat onmiddellijk de intuïtie begeleidt), omdat zijn bepalende grond kan liggen in het concept van datgene wat als het bovenzintuiglijke substratum van de menselijkheid kan worden beschouwd.

De oplossing van een antinomie bestaat slechts hierin dat twee ogenschijnlijk strijdige stellingen niet werkelijk met elkaar in tegenspraak zijn, maar met elkaar verenigbaar kunnen zijn, ook al overstijgt de verklaring van de mogelijkheid van hun concept ons kenningsvermogen. Dat deze schijn ook natuurlijk en onvermijdelijk is voor de menselijke rede, en dus waarom hij bestaat en blijft voortbestaan, ook al bedriegt hij niet meer na de oplossing van de schijnbare tegenstrijdigheid, kan ook op deze grond begrijpelijk worden gemaakt.

Wij nemen namelijk het concept, waarop de universele geldigheid van een oordeel moet berusten, in dezelfde zin in beide strijdige stellingen, en toch beweren wij twee tegengestelde predicaten ervan. Zo zou de these moeten luiden: het smaakoordeel berust niet op \textbf{bepaalde} concepten; maar in de antithese zou het moeten luiden: het smaakoordeel berust toch op enig, hoewel \textbf{onbepaald} concept (namelijk van het bovenzintuiglijke substratum der verschijningen); en dan zou er geen strijd tussen hen zijn.

Wij kunnen niet meer doen dan deze tegenstrijdigheid in de eisen en tegen-eisen van de smaak wegnemen. Een bepaald objectief principe van de smaak aan te geven, waardoor haar oordelen geleid, onderzocht en bewezen zouden kunnen worden, is volkomen onmogelijk; want dan zou het geen smaakoordeel zijn. Het subjectieve principe, namelijk de onbepaalde idee van het bovenzintuiglijke in ons, kan slechts worden aangewezen als de enige sleutel om dit vermogen, dat zelfs in zijn bronnen voor ons verborgen is, te ontrafelen, maar er is niets waardoor het verder begrijpelijk gemaakt kan worden.

De hier uiteengezette en opgeloste antinomie berust op het juiste begrip van de smaak, namelijk als een louter reflecterend vermogen van de esthetische oordeelskracht; en de twee ogenschijnlijk strijdige grondstellingen zouden hiermee met elkaar verenigd kunnen worden, voor zover \textbf{beide waar kunnen zijn}, wat ook voldoende is. Als men daartegen het \textbf{aangename} als de bepalende grond van de smaak zou aannemen (wegens de enkelvoudigheid van de voorstelling die ten grondslag ligt aan het smaakoordeel), zoals sommigen doen, of het principe van \textbf{volmaaktheid}, zoals anderen doen (wegens de universele geldigheid ervan), en het smaakoordeel dienovereenkomstig zou vaststellen, dan zou daaruit een antinomie voortvloeien die helemaal niet opgelost kan worden, tenzij men aantoont dat \textbf{beide} tegenovergestelde (maar niet slechts tegengestelde) \textbf{stellingen onwaar zijn}: wat dan zou bewijzen dat het concept waarop elk berust in zichzelf tegenstrijdig is. Zo ziet men dat de opheffing van de antinomie van de esthetische oordeelskracht eenzelfde weg volgt als die welke de \textit{Kritiek} inslaat bij de oplossing van de antinomieën van de zuivere theoretische rede, en dat men hier, evenals in de \textit{Kritiek van de praktische rede}, tegen zijn wil gedwongen wordt om voorbij het zintuiglijke te kijken en het verenigende punt van al onze vermogens \textit{a priori} in het bovenzintuiglijke te zoeken: omdat geen andere weg overblijft om de rede met zichzelf in overeenstemming te brengen.

\begin{mdframed}[leftmargin=0pt,rightmargin=0pt,backgroundcolor=blue!6,linecolor=white,innertopmargin=6pt,innerbottommargin=6pt]

De schijnbare tegenstrijdigheid die ontstaat wanneer wij beweren dat smaakoordelen enerzijds aanspraak maken op universele geldigheid, maar anderzijds niet op bepaalde, determineerbare begrippen berusten, is geen echte tegenspraak, maar een natuurlijke illusie die voortvloeit uit een misverstand over de aard der begrippen waarop deze oordelen zich baseren. Het smaakoordeel is niet zonder grond; het veronderstelt weliswaar een concept, maar niet in de zin van een empirisch of logisch vastomlijnd begrip, zoals het verstand die gebruikt om kennis te verwerven. Nee, het is een begrip van een geheel andere orde: het transcendentaal idee van het bovenzintuiglijke, dat als grond dient voor alle zintuiglijke ervaring, zonder zelf ooit in de zintuigen te kunnen worden gevangen of theoretisch te kunnen worden vastgelegd.

Wanneer wij oordelen dat iets mooi is, doen wij dat niet louter op basis van persoonlijke voorkeur, alsof ons oordeel niets meer is dan een uitdrukking van genot of afkeer. Een dergelijk oordeel zou indrukwekkend subjectief zijn, beperkt tot het individu: \textit{Dit bevalt mij, maar een ander kan het anders ervaren.} Toch maken wij met elk smaakoordeel aanspraak op instemming, alsof schoonheid niet alleen voor ons, maar voor iedereen geldt. Deze aanspraak kan alleen gerechtvaardigd zijn als het smaakoordeel gerelateerd is aan een grond die verder reikt dan louter subjectieve sensatie. Die grond is geen gewoon begrip, geen vastomlijnde definitie die wij kunnen uitleggen of bewijzen, maar het zuivere, rationele concept van het bovenzintuiglijke; een idee dat het voorwerp (en het oordelende subject zelf) niet als ding-op-zich beschouwt, maar als verschijning, als iets dat alleen in relatie tot ons kennend vermogen bestaat.

Het bovenzintuiglijke is geen voorwerp van kennis; het is niet iets dat wij kunnen vastpinnen in woorden of analyseren met het verstand. Toch is het de onmisbare voorwaarde waarbinnen smaak zijn universele aanspraak kan doen gelden. Zonder deze veronderstelling zou het smaakoordeel geen recht hebben op de instemming van anderen; het zou slechts een willekeurige, persoonlijke reactie zijn, zonder enige claim op algemene betekenis. De schijn van tegenstrijdigheid ontstaat doordat wij geneigd zijn om het begrip waarop het smaakoordeel berust op dezelfde manier te begrijpen als de begrippen die wij in de wetenschap of de logica gebruiken: als iets dat duidelijk kan worden gedefinieerd en waaruit conclusies kunnen worden getrokken. Maar het bovenzintuiglijke is onbepaalbaar; het is geen object van kennis, maar een regulatief principe, een idee dat ons denkvormend vermogen stuurt zonder zelf ooit volledig kenbaar te zijn.

Daarom kan het smaakoordeel nooit worden ``bewezen'' in de gebruikelijke zin van het woord. Het is geen kwestie van redeneren of argumenteren alsof het om een wetenschappelijke stelling gaat. Toch is het ook niet willekeurig, want het berust op een subjectieve, maar universele structuur van onze geest: de idee dat de natuur, in haar verschijningen, op een bepaalde manier doelmatig is voor ons oordeelsvermogen, zonder dat wij kunnen aangeven hoe dat precies werkt. De oplossing van deze antinomie ligt in het inzien dat de twee ogenschijnlijk strijdige stellingen (\textit{smaak berust niet op begrippen} en \textit{smaak berust wel op een begrip}) elkaar niet werkelijk tegenspreken, maar naar verschillende aspecten van hetzelfde vermogen verwijzen. De eerste stelling benadrukt dat smaak niet berust op determineerbare, empirische of logische begrippen; de tweede dat het wel degelijk een grond heeft in een idee dat, hoewel onkenbaar, noodzakelijk is om de universele aanspraak van het smaakoordeel te rechtvaardigen.

Dit idee is het bovenzintuiglijke, dat als een substratum fungeert voor zowel het voorwerp als het subject, en zo de brug slaat tussen onze individuele ervaring en de eis van algemene geldigheid. Wanneer men daartegen zou proberen om smaak te gronden in louter genot (alsof schoonheid niets meer is dan aangename sensatie) of in een begrip van volmaaktheid (alsof schoonheid een meetbare eigenschap is), dan zou men in een echte, onoplosbare tegenstrijdigheid vervallen. Beide benaderingen mislukken, omdat zij het wezenlijke karakter van het smaakoordeel negeren: dat het noch louter subjectief, noch objectief in de zin van wetenschappelijke kennis is. Het is juist de reflecterende aard van de smaak, het vermogen om zonder vaste regels te oordelen, maar toch met een claim op universaliteit, die ons dwingt om voorbij het zintuiglijke te kijken.

Net zoals in de theoretische filosofie, waar de rede gedwongen wordt om ideeën als die van de ziel, de wereld als geheel, en God te denken om haar eigen eenheid te bewaren, zo moet ook de esthetische oordeelskracht een beroep doen op het bovenzintuiglijke om haar eigen wetmatigheid te rechtvaardigen. Dit inzicht toont aan dat de menselijke geest niet beperkt kan blijven tot de wereld van de zintuigen en het verstand, maar altijd al verwijst naar iets dat daarbuiten ligt; niet als een voorwerp van kennis, maar als een noodzakelijke vooronderstelling voor de betekenis die wij aan de wereld geven. De smaak, in al haar subjectiviteit, openbaart zo een diepere laag van onze ervaring: de erkenning dat wij niet alleen wezens zijn die waarnemen en redeneren, maar ook wezens die zin zoeken in wat zij ervaren, en die daartoe een beroep doen op ideeën die ons overstijgen. In die zin is de antinomie van de smaak geen probleem, maar een openbaring: een teken dat onze ervaring van schoonheid ons altijd al verbindt met iets dat groter is dan onszelf.

\bigskip
Deze paragraaf onthult de diepste structuur van Kants esthetica en tegelijkertijd een fundamentele waarheid over de menselijke geest: schoonheid is de brug tussen onze individuele ervaring en een universele orde die wij niet kunnen kennen, maar wel moeten veronderstellen. Wat op het eerste gezicht een technisch-filosofische oplossing lijkt voor een logische tegenstrijdigheid, blijkt in werkelijkheid een openbaring over hoe wij betekenis geven aan de wereld. De antinomie van de smaak is geen probleem dat opgelost moet worden, maar een spiegel die ons laat zien wie wij zijn als denkende en voelende wezens.

Allereerst toont Kant hier aan dat ware schoonheid niet kan worden ingesloten in begrippen, noch kan worden gereduceerd tot louter subjectieve sensatie. Het smaakoordeel is geen kwestie van persoonlijke smaak, maar ook geen objectief feit dat bewezen kan worden. Er is een derde weg: een ervaring die ons zowel als individuen als als gemeenschap raakt. Dit is geen zwakte, maar juist de kracht van esthetische oordelen. Schoonheid is geen zaak van weten, maar van ervaren en delen: een activiteit die ons dwingt om voorbij onze eigen subjectiviteit te kijken, zonder ons te verlaten op vaste, objectieve criteria. Voor vrijmetselaren is dit een herkenbare gedachte: ook symbolen en rituelen functioneren op deze manier. Zij hebben geen eenduidige betekenis die kan worden uitgelegd of bewezen, en toch verbinden zij diegene die ermee werken. Net zoals schoonheid een beroep doet op een onbepaalbaar, maar noodzakelijk idee (het bovenzintuiglijke), zo verwijzen maçonnieke symbolen naar een diepere laag van betekenis die nooit volledig in woorden kan worden gevangen, maar die wel ervaren en gedeeld kan worden.

Wat Kant hier blootlegt, is niets minder dan de structuur van menselijke zingeving zelf. Wij zijn wezens die niet alleen waarnemen en redeneren, maar ook zin zoeken in wat wij ervaren. Het bovenzintuiglijke is geen mystiek concept, maar de erkenning dat onze ervaring van de wereld altijd al verwijst naar iets dat voorbij het zintuiglijke ligt. Dit is geen vlucht uit de werkelijkheid, maar juist de erkenning dat de werkelijkheid zelf open is, dat zij meer is dan wat wij kunnen meten of definiëren. Het smaakoordeel is daarom geen louter esthetische kwestie, maar een fundamentele houding ten opzichte van de wereld: de bereidheid om te ervaren wat niet volledig kan worden begrepen, en toch te geloven dat het betekenis heeft.

Deze paragraaf is ook een kritiek op twee uitersten die de moderne mens vaak in vervoering brengen: enerzijds het reductionisme, dat alles wil terugbrengen tot meetbare feiten en persoonlijke voorkeuren, en anderzijds het dogmatisme, dat absolute waarheden claimt waar geen ruimte is voor twijfel of interpretatie. Kant laat zien dat ware wijsheid niet in een van beide uitersten ligt, maar in het vermogen om te leven met de spanning tussen subjectiviteit en universaliteit, tussen ervaring en idee. Dit is precies wat het maçonnieke pad ons leert: ware kennis is geen kwestie van bezit, maar van zoektocht, niet van vaste antwoorden, maar van open vragen die ons verbinden met anderen en met iets dat groter is dan onszelf.

Bovendien onthult deze paragraaf dat vrijheid en verbondenheid geen tegengestelden zijn, maar twee kanten van dezelfde ervaring. Het smaakoordeel is vrij, omdat het niet gedwongen kan worden; het is verbindend, omdat het een beroep doet op een gemeenschappelijke grond. Dit is de kern van het maçonnieke ideaal van broederschap: wij zijn vrij om onze eigen weg te gaan, maar die vrijheid vindt haar betekenis alleen in de verbondenheid met anderen die dezelfde zoektocht delen. Het bovenzintuiglijke is dus geen abstract concept, maar de levende werkelijkheid van een gemeenschap die samen zoekt naar licht.

Ten slotte is deze paragraaf een uitnodiging tot nederigheid en verwondering. Kant leert ons dat de diepste waarheden niet kunnen worden bewezen of uitgelegd, maar alleen ervaren en gedeeld. Dit is geen beperking, maar een bevrijding: het inzicht dat wij niet alles hoeven te weten om betekenis te vinden. Schoonheid, net als wijsheid, is geen bestemming, maar een manier van onderweg zijn, een weg die wij alleen kunnen gaan door ons open te stellen voor de ervaring van anderen en voor de mysteries die ons overstijgen. In die zin is de oplossing van de antinomie van de smaak niet alleen een filosofisch inzicht, maar een levenshouding: de bereidheid om te leven met wat wij niet volledig kunnen begrijpen, en toch te geloven dat het de moeite waard is om te zoeken.

\end{mdframed}
